% Source: https://artofproblemsolving.com/community/c3h2631445_complete_factorization_problem
% Stable identity: aops-topic-2631445-q01
Since $\left(x^6-1\right)\left(x^{12}+x^6+1\right)=x^{18}-1$, the roots of $x^{12}+x^6+1$ are the $18$th roots of unity that are not $6$th roots of unity:
\[z=\cos\theta+i\sin\theta,\qquad\theta=20^\circ k,\quad k\in\{1,2,4,5,7,8,10,11,13,14,16,17\}.\]
None is real ($\theta\ne0^\circ,180^\circ$), so there are no linear factors. Pairing each root with its conjugate ($\theta$ with $360^\circ-\theta$) gives the real quadratic $(x-z)(x-\bar z)=x^2-2\cos\theta\,x+1$, which has no real roots because its discriminant $4\cos^2\theta-4$ is negative. With $\theta=20^\circ,40^\circ,80^\circ,100^\circ,140^\circ,160^\circ$ and $\cos(180^\circ-\theta)=-\cos\theta$:
\[\boxed{\begin{aligned}x^{12}+x^6+1={}&\left(x^2-2\cos20^\circ\,x+1\right)\left(x^2+2\cos20^\circ\,x+1\right)\left(x^2-2\cos40^\circ\,x+1\right)\\&\cdot\left(x^2+2\cos40^\circ\,x+1\right)\left(x^2-2\cos80^\circ\,x+1\right)\left(x^2+2\cos80^\circ\,x+1\right)\end{aligned}}\]
Over the rationals the factorisation stops earlier: $x^{12}+x^6+1=\left(x^6+1\right)^2-x^6=\left(x^6+x^3+1\right)\left(x^6-x^3+1\right)$.
